% =============================================================================
%  Methods table: how each pipeline step contributes to the HTML dashboard
%  Compile standalone (pdflatex methods_table.tex) or copy the longtable
%  (from \begin{longtable} to \end{longtable}) into the manuscript.
%  Needs: booktabs, longtable, array, ragged2e (all in TeX Live).
% =============================================================================
\documentclass[a4paper,10pt]{article}
\usepackage[margin=2cm]{geometry}
\usepackage[T1]{fontenc}
\usepackage[scaled=0.92]{helvet}
\renewcommand{\familydefault}{\sfdefault}
\usepackage{booktabs,longtable,array,ragged2e}

\newcolumntype{L}[1]{>{\RaggedRight\arraybackslash}p{#1}}
\newcommand{\file}[1]{\texttt{\detokenize{#1}}}   % file names with underscores

\begin{document}
\pagestyle{empty}
\footnotesize
\setlength{\tabcolsep}{4pt}
\renewcommand{\arraystretch}{1.18}

\begin{longtable}{@{}L{2.3cm} L{3.2cm} L{5.9cm} L{4.4cm}@{}}
\caption{Pipeline steps and their contribution to the interactive HTML dashboard.
Numbers refer to the current run (14 organs, 91 organ pairs); papers are unique papers unless stated otherwise.}
\label{tab:html-generation}\\
\toprule
\textbf{Step} & \textbf{Input} & \textbf{Processing} & \textbf{Use in the HTML dashboard}\\
\midrule
\endfirsthead
\toprule
\textbf{Step} & \textbf{Input} & \textbf{Processing} & \textbf{Use in the HTML dashboard}\\
\midrule
\endhead
\midrule \multicolumn{4}{r}{\textit{continued on next page}}\\
\endfoot
\bottomrule
\endlastfoot

\textbf{1}~Literature search and filtering
 & 14 organs, 91 organ pairs; organ MeSH terms and aliases; 22 metabolic and 21 crosstalk keywords
 & One PubMed query per pair (title/abstract), last 20 years, up to 20\,000 hits, 0.4~s between requests;
   PMID, title, abstract, year and MeSH descriptors fetched via EFetch.
   Local rule-based filters keep a paper only if both organs and a crosstalk keyword occur in the
   same sentence (whole-word, negation-aware) and no third organ indicates a different relationship
   (multi-organ disambiguation).
   187\,757 hits retrieved (summed over pairs); 27\,416 pair--paper assignments kept, i.e.\ 22\,451 unique
   papers (87 of 91 pairs with at least one paper).
 & Existence of a connection and its paper count; overlap-coefficient (OO) based edge weight;
   paper list with abstracts per connection; Literature Statistics tab (Sankey and bar plots).\\[3pt]

\textbf{2}~Condition subselection
 & Papers of step 1; 7 healthy and 11 obese keywords
 & A paper is kept for a condition if its title or abstract contains at least one keyword of that condition
   (whole-word, negation-aware, anywhere in the text). The two sets overlap (200 papers).
   Healthy: 1\,585 papers in 70 pairs; obese: 3\,594 papers in 72 pairs.
   Computed on the fly while the HTML is built.
 & Healthy/Obese switch: a connection is shown if it has at least one matching paper; paper counts,
   key players, OO weights and statistics are recomputed from the filtered papers.\\[3pt]

\textbf{3}~Key player extraction
 & Title, abstract and MeSH descriptors of the papers of each pair and condition
 & Multi-pattern (Aho--Corasick) whole-word matching against curated vocabularies of hormones (171),
   metabolites (254) and proteins/transporters (45\,138). Synonyms and paralog families are merged to one
   display name at match time; each term counts once per paper (count = number of papers).
 & Ranked lists of hormones, metabolites and proteins with paper counts in the information panel of each
   connection; recalculated for the Healthy and Obese views.\\[3pt]

\textbf{4}~Connection-type classification (LLM)
 & First 100 papers of a pair (abstracts truncated to 400 characters); definitions of 11 connection types
 & Local LLM (llama3.2 via Ollama, temperature 0.5) assigns 1--3 types per run; 5 independent runs;
   types named in the majority of runs (at least 3 of 5) are kept; if none reaches a majority the single
   most frequent type is kept.
 & Connection-type label and colour of each edge and the type legend; Healthy/Obese views use the
   condition-specific classification where it exists and the all-papers classification otherwise.\\[3pt]

\textbf{5}~Connection descriptions (LLM)
 & Papers of a pair (condition-filtered for the condition files)
 & 25 papers drawn at random (all if fewer); the LLM selects the 5 most relevant; a second prompt writes a
   3--4 sentence summary of the metabolic and hormonal basis of the axis that cites every selected paper.
 & Cited summary with the five source papers in the information panel of each connection.\\[3pt]

\textbf{6}~Organ descriptions (LLM)
 & 14 organs; organ MeSH term/aliases and the 22 metabolic keywords
 & Separate PubMed search per organ (title/abstract, last 5 years, up to 50 papers; 700 in total); the first six
   papers (abstracts truncated to 400 characters) are given to the LLM, which writes exactly five sentences on
   consumed and produced substrates, energy pathways and hormonal regulation. Refusals are detected and
   re-prompted with a simplified prompt.
 & Description and source papers in the information panel of each organ (node).\\[3pt]

\textbf{7}~Assembly of the dashboard
 & Caches of steps 1, 4, 5 and 6 (plus the condition-specific files); PET partial-correlation networks
   of the healthy and obese cohorts
 & Undirected NetworkX graph (14 nodes, 91 edges) with all edge and node attributes, including the
   precomputed condition variants; export to a single HTML file rendered with Cytoscape.js (version 3.29.2,
   loaded from a CDN). The cohort partial-correlation networks (16 healthy and 9 obese edges) are added
   as comparison networks.
 & Interactive network with All/Healthy/Obese views, comparison with the PET-based networks, and the
   Literature and Edge Statistics tabs.\\

\end{longtable}
\end{document}
